\section{Conclusion}\label{sec:conclusion}

When downstream activity moves away from a node that produces information, the consequences are not determined solely by where sales occur. The node that incurs the cost of generating information may lose the transactions over which it appropriates a return, while the information itself can remain valuable after being incorporated into an upstream product used through other routes.

This paper formalizes that tension as information appropriability under downstream reallocation. In the general model, a shrinking appropriable footprint reduces private information effort. Coordinated information effort instead rises when the weighted marginal value created through the expanding route exceeds the value lost through the shrinking route. Under that condition, downstream reallocation generates the central sign divergence
\[
\frac{d\Eprivate}{d\eta}<0<\frac{d\Ecoordinated}{d\eta}.
\]
Cross-route reuse is essential: when the expanding route receives no product-improvement benefit, the positive coordinated response disappears.

A microfounded benchmark further shows that private effort is below coordinated effort, coordinated effort is below socially optimal effort, and the private--system wedges widen with reallocation under the stated restrictions. Around a valid private retention threshold, the same logic can produce a local region in which private organization drops an information-producing downstream capacity even though social welfare favors retention.

The broader organizational lesson is limited but consequential. Reallocating transactions can weaken the private incentive to produce knowledge even when the economic value of that knowledge is increasing. Whether this tension matters in a particular market depends on how costly information production is, how much of its return the originating node can appropriate, and how reusable the resulting product improvement is across downstream routes.
